\section{Verification and Validation}
\label{sec:validation}

\subsection{Strategy}

The application makes a falsifiable claim --- that its output is a physically
motivated simulation rather than an aesthetic approximation --- and the test
strategy is built to falsify it. Four layers:

\begin{enumerate}
\item \textbf{Unit tests} over \texttt{EmulsionCore}, which contains the
mathematics and no GPU code. Fast, hermetic, run on every commit.
\item \textbf{Reference implementation cross-check.} Every shader has a
corresponding CPU implementation in Swift. Random inputs are pushed through both
and compared. This catches the entire class of GPU-specific defects ---
precision, edge clamping, indexing --- that visual inspection misses.
\item \textbf{Golden image tests.} A corpus of source images rendered with a
corpus of recipes, compared against stored references with a perceptual metric.
\item \textbf{Sensitometric validation.} Synthetic step wedges transported
through the pipeline and measured, compared against published and measured film
data. This is what tests the claim itself.
\end{enumerate}

\subsection{Unit Test Coverage of the Model}

\begin{table}[htbp]
\caption{Model Unit Tests}
\label{tab:unittests}
\begin{center}
\renewcommand{\arraystretch}{1.12}
\begin{tabularx}{\columnwidth}{@{}l>{\raggedright\arraybackslash}X@{}}
\toprule
\textbf{ID} & \textbf{Assertion} \\
\midrule
V-01 & $\softp{a}{u} \to \max(u,0)$ as $a \to 0$; relative error $< 10^{-6}$ at $a = 10^{-4}$. \\
V-02 & Stable and naive softplus agree to $10^{-6}$ for $|u/a| < 30$; stable form finite for $|u/a| < 10^{6}$. \\
V-03 & \eqref{eq:charcurve} straight-line slope equals $\gamma$ to $10^{-4}$ at the curve midpoint. \\
V-04 & \eqref{eq:charcurve} attains $\Dmax$ within $0.02$ when \eqref{eq:wellformed} holds, for 10{,}000 random valid parameter sets. \\
V-05 & \eqref{eq:pointgamma} matches a central-difference derivative of \eqref{eq:charcurve} to $10^{-5}$. \\
V-06 & Monotonicity: $D(x_{i+1}) \geq D(x_i)$ over a $10^4$-point sweep, for every reachable UI state (property test). \\
V-07 & \eqref{eq:speedpoint} inverts \eqref{eq:charcurve}: $D(x_{\mathrm{sp}}) = \Dmin + 0.10 \pm 10^{-4}$. \\
V-08 & Aim balance \eqref{eq:aimbalance} yields neutral print for every (negative, print) pair: $\Delta E_{00} < 0.5$. \\
V-09 & Halation PSF integrates to unity: $\left|\int h - 1\right| < 10^{-3}$ under discretization. \\
V-10 & Gaussian mixture \eqref{eq:gaussmix} matches \eqref{eq:halotf} to $< 3\%$ relative over $[0, 1/2\ell]$. \\
V-11 & Grain variance at half resolution equals the Selwyn prediction from full resolution to $10\%$. \\
V-12 & Grain field is bit-identical across two runs with the same seed. \\
V-13 & Grain field is identical at the same film coordinate for render extents differing by $4\times$. \\
V-14 & Interlayer operator is mean-preserving: uniform input is unchanged to $10^{-6}$. \\
V-15 & Tetrahedral interpolation reproduces neutrals exactly: $|R - G|, |G - B| < 10^{-6}$ for neutral input. \\
V-16 & Null test (AC-7): identity configuration matches direct color conversion, $\Delta E_{00} < 0.5$. \\
V-17 & Recipe JSON round-trips through every migration version to an identical resolved parameter set. \\
V-18 & Graph aliasing: no two live-overlapping resources share a heap offset (property test, $10^4$ random graphs). \\
V-19 & Graph dirty propagation matches the specification table (generates Table~\ref{tab:dirtycost}). \\
V-20 & Dead pass elimination removes exactly the expected passes for each bypass combination. \\
\bottomrule
\end{tabularx}
\end{center}
\end{table}

\subsection{Sensitometric Validation}

This is the acceptance test for the central claim, and it is run as a separate
suite because it requires reference data.

\paragraph{Method} A synthetic 21-step wedge spanning $4.0$ in $\logE$ is
constructed in the working space, transported through the negative model, and
the resulting densities read. The measured curve is compared point-by-point
against the reference $D$--$\logE$ data digitized from the manufacturer's
datasheet \cite{kodakportra,kodakektar,kodakvision3,kodaktrix,fujidata,ilfordhp5}
and, where available, against densitometer readings of a physically exposed and
processed step wedge measured under ISO~5-3 conditions \cite{iso5}.

\paragraph{Metric} Maximum absolute density error over the range $\Dmin + 0.1$
to $\Dmax - 0.1$, per record. AC-1 tolerance is $0.03$.

\paragraph{Result format} Each stock reports a table of $(\Delta_R, \Delta_G,
\Delta_B)$ maximum errors and an RMS error. A stock failing AC-1 does not ship;
it returns to the profile fitter.

The same procedure is applied end-to-end for AC-2: the wedge is transported
through both negative and print, and the resulting print densities are compared
against a reference print-through of the same wedge.

\subsection{Grain Spectrum Validation}

Grain is validated in the frequency domain, which is the only way to test it
meaningfully.

A uniform mid-density field is rendered with grain enabled. The output is
converted to density, the mean removed, and a periodogram computed with a Hann
window and radial averaging. The result is compared against
\eqref{eq:wienerclosed} evaluated with the profile's parameters, and against
published Wiener spectra where they exist.

Two failure modes this catches that visual inspection does not:

\begin{itemize}
\item \textbf{Lattice aliasing.} If the grain lattice pitch is close to the
pixel pitch, the spectrum shows a spike at the lattice frequency. Visually this
reads as a faint regular texture that most reviewers will describe as "fine" but
which is plainly wrong in the spectrum.
\item \textbf{Insufficient decorrelation.} If the hash function has poor
avalanche behavior in one bit position, adjacent lattice cells correlate and the
spectrum shows directional structure. This is invisible at normal viewing
distance and obvious at 400\%.
\end{itemize}

RMS granularity is validated separately by measuring $\sigma_D$ through a
simulated $48\,\mu$m aperture across a density sweep and comparing the resulting
curve against \eqref{eq:grainvar} and against published RMS granularity figures
(AC-3). The equivalent measurement on the digital output path follows the noise
methodology of ISO~15739 \cite{iso15739}, which is how the grain figure is
compared against sensor noise rather than confused with it.

\subsection{Halation Validation}

A synthetic scene containing a small disc of known radiance on a black field is
rendered. The radial profile of the resulting halo is extracted by azimuthal
averaging, the direct component subtracted, and an exponential fitted. The fitted
decay length is compared against $\ell_c$ (AC-4), and the ratio of peak halo
intensities across channels against $\boldsymbol{\beta}$.

Additionally, the energy conservation property of \eqref{eq:haladd} is tested:
total scene energy before and after the halation stage must agree to $0.5\%$
when the source term equals the full image.

\subsection{Calibration}
\label{sec:calibration}

The constant $g_{\mathrm{cal}}$ of \eqref{eq:anchor} relates the decoded
middle-gray value to the working-space unit, and it must be determined per
device generation. The procedure:

\begin{enumerate}
\item Photograph a calibrated 18\% reflectance gray card under diffuse
illumination, at a metered exposure, with a fixed manual setting.
\item Decode with the configuration of Listing~\ref{lst:rawconfig} and convert
to ACEScg.
\item Measure the mean of a central patch. This value is
$g_{\mathrm{cal}}$.
\item Repeat across illuminants (D65, D50, tungsten, fluorescent) and confirm
the value is stable to within $\pm 3\%$; a larger spread indicates the white
balance path is contaminating the anchor.
\end{enumerate}

Measured $g_{\mathrm{cal}}$ values are stored in a device table keyed by
\texttt{utsname.machine}, with a fallback derived from the device's reported
color matrix for unknown devices. The fallback's accuracy is the reason unknown
new devices are flagged in telemetry.

\subsection{Perceptual Evaluation}

Numerical validation establishes that the model is correct. It does not
establish that the result is good, and the two are genuinely different
questions.

A panel of twelve evaluators --- six matching persona P1 (working
photographers who shoot film) and six matching P2 (colorists) --- performs a
two-alternative forced choice task. For each of forty scenes, the panel is shown
the \EM{} render alongside a lab scan of the same scene shot on the
corresponding real stock, in randomized order, and asked which is the
photographic film. The target is a discrimination rate below 65\%; chance is
50\%.

This test is deliberately harsh, and failing it is expected in early rounds. Its
value is diagnostic: the panel is asked, after each choice, what cue they used.
The accumulated cue list is the prioritized defect list for the imaging team,
and in early prototyping it is a far more efficient signal than open-ended
critique.

A secondary evaluation asks the panel to rank \EM{} against the leading desktop
film emulation tools on the same source material, which establishes competitive
position rather than absolute fidelity.

\subsection{Device and Configuration Matrix}

\begin{table}[htbp]
\caption{Test Matrix}
\label{tab:testmatrix}
\begin{center}
\renewcommand{\arraystretch}{1.15}
\begin{tabularx}{\columnwidth}{@{}l>{\raggedright\arraybackslash}X@{}}
\toprule
\textbf{Dimension} & \textbf{Values} \\
\midrule
Device & iPhone 12 (A14), 13 Pro (A15), 15 Pro (A17 Pro), 16 Pro (A18 Pro), iPad Pro M4 \\
OS & Minimum supported, current release, current beta \\
Capture & ProRAW 12\,MP, ProRAW 48\,MP, Bayer RAW, imported DNG, imported TIFF \\
Appearance & Light, dark, increased contrast, reduce transparency \\
Dynamic Type & Default, XXL \\
Locale & en-US, de-DE (long strings), ja-JP (CJK), ar-EG (RTL) \\
Storage & Ample, near-full (cache eviction path) \\
Thermal & Nominal, artificially induced serious \\
\bottomrule
\end{tabularx}
\end{center}
\end{table}

\subsection{Continuous Integration}

Every pull request runs: build for all configurations; \texttt{EmulsionCore}
unit tests (target: under 30\,s); reference cross-check on the simulator;
golden image tests on a physically attached A17 device in the build farm; a
static analysis pass; and a performance smoke test that fails the build if the
preview render regresses more than 10\% against the recorded baseline.

Golden image comparison uses $\Delta E_{00}$ with a mask excluding the grain
contribution --- computed by rendering the same recipe with $g = 0$ and
comparing that --- since grain is stochastic across code changes that alter
dispatch order even when statistically identical. The grain itself is validated
by the spectral test rather than by image comparison, which is the correct
division: deterministic parts compared by image, stochastic parts compared by
statistic.

\subsection{Known Limits of the Model}

Honest specification requires stating what the model does not capture. These are
accepted limitations for v1.0, documented so that they are not rediscovered as
bugs:

\begin{itemize}
\item \textbf{Reciprocity failure} is not modelled. Very long exposures cause
nonlinear speed loss that differs per layer. The parameter hooks exist
(\eqref{eq:tungsten} accepts a per-layer offset) but no exposure-time-dependent
model is fitted. This matters for simulated long exposures only.
\item \textbf{True spectral rendering} is absent; the reduced spectral strategy
bakes crosstalk at profile time against a reference illuminant. Scenes with
strongly non-blackbody illuminants (sodium vapor, some LEDs) will render with
crosstalk that is subtly wrong.
\item \textbf{Emulsion MTF} is modelled only through the grain kernel and the
interlayer effect; the light-scattering resolution loss within the emulsion is
not separately modelled. The visible consequence is that \EM{} renders slightly
sharper than the corresponding real stock at high enlargement.
\item \textbf{Scanner characteristics} are not modelled. Most users' reference
for "what film looks like" is a scan, and scanners impose their own transfer,
sharpening, and color handling. The "bypass" print profile partially addresses
this, but a proper scanner model is future work.
\end{itemize}
